\documentclass[12pt]{article}

\usepackage[margin=1in]{geometry}
\usepackage{setspace}
\usepackage{parskip}
\usepackage{titlesec}
\usepackage{hyperref}
\usepackage{booktabs}
\usepackage{amsmath}

\hypersetup{colorlinks=true, linkcolor=blue, citecolor=blue, urlcolor=blue}

\titleformat{\section}{\large\bfseries}{}{0em}{}
\titlespacing{\section}{0pt}{14pt}{4pt}

\setlength{\parskip}{6pt}
\onehalfspacing

\title{\textbf{CNN-Based Re-ranking for Image Similarity Retrieval}\\[0.4em]
\large Deep Learning Course Project Proposal}
\author{[Author(s)]}
\date{\today}

\begin{document}
\maketitle

% ---------------------------------------------------------------
\section{Background and Motivation}

Image similarity retrieval---given a query image, return the most visually similar images from a large gallery---underlies a wide range of applications including visual product search, medical image analysis, and content-based recommendation. The prevailing approach encodes images into dense embedding vectors using a pre-trained deep network, then ranks gallery items by cosine similarity to the query embedding. This pipeline is efficient and generalizes well, but it has a notable limitation: cosine similarity is a fixed, non-adaptive metric. It treats every dimension of the embedding space equally and has no mechanism to learn what ``similar'' should mean for a specific domain or task.

A natural remedy is to add a re-ranking stage after the initial retrieval. Instead of returning the cosine-ranked list directly, one applies a second, more powerful model to re-score the top candidates. Prior work has demonstrated clear gains from this idea---most notably, transformer-based re-rankers that attend jointly over query--candidate pairs. However, transformers require substantial compute and are hard to train from scratch with limited data and GPU resources, which is a real constraint in a course project setting. This raises a practical question: can a much simpler CNN-based re-ranker, trained on pairwise labels derived from category annotations, still produce meaningful improvements over cosine similarity while remaining tractable? That is the question we set out to answer.

% ---------------------------------------------------------------
\section{Research Objectives}

This project has two closely related goals. The primary goal is to design and evaluate a lightweight CNN that takes a query--candidate image pair as input and produces a refined similarity score, replacing the cosine similarity used in the final ranking step. The secondary goal is to understand where this re-ranker helps and where it does not---specifically, how the choice of backbone encoder and the depth of the CNN affect the quality--efficiency trade-off. Together, these goals will give us a clearer picture of how much headroom exists above the cosine similarity baseline, and how much model complexity is actually needed to capture it.

% ---------------------------------------------------------------
\section{Proposed Methodology}

\paragraph{Overall pipeline.}
The system is structured as a two-stage retrieval pipeline. In the first stage, a pre-trained encoder (we will experiment with CLIP, DINOv2, ResNet, and EfficientNet) extracts a feature embedding for every image, and the Top-$K$ gallery candidates for a given query are retrieved by cosine similarity. This coarse stage is fast and sets the recall ceiling. In the second stage, each of the $K$ candidates is re-scored by a learned pairwise model, and the final ranked list is produced by sorting on the new scores.

\paragraph{CNN re-ranker design.}
For the re-ranker, we concatenate the query image and the candidate image along the channel dimension to form a $6 \times H \times W$ input tensor. This gives the network direct pixel-level access to both images simultaneously, allowing it to detect fine-grained differences that are invisible in the aggregated embedding space. The network itself is kept deliberately simple: a stack of Conv--ReLU--MaxPool blocks for feature extraction, followed by a small fully connected head that outputs a single similarity score $\hat{s} \in [0, 1]$ via sigmoid activation.

\paragraph{Training.}
The re-ranker is trained with binary cross-entropy. Positive pairs are formed from images of the same category; negative pairs come from different categories, with hard negatives sampled from the top of the cosine-ranked list (i.e., candidates that the baseline retrieves but are actually incorrect). Focusing on hard negatives is important because it forces the model to learn distinctions that cosine similarity already struggles with, rather than simply separating easy cases that would be correctly ranked anyway.

% ---------------------------------------------------------------
\section{Evaluation Plan}

We will report Recall@$K$ and Top-$K$ Accuracy for $K \in \{1, 5, 10\}$, comparing three conditions: (1) a backbone with cosine similarity and no re-ranking, (2) the same backbone plus our CNN re-ranker, and (3) the same setup across different backbone choices. Beyond aggregate metrics, we will include qualitative side-by-side comparisons of baseline and re-ranked results for selected query images, which often reveal failure modes that numbers alone do not capture. For ablations, we plan to vary the backbone, the value of $K$, and the depth of the CNN to understand sensitivity to each design choice.

% ---------------------------------------------------------------
\section{Timeline}

\begin{tabular}{cl}
\toprule
\textbf{Week} & \textbf{Task} \\
\midrule
1--2 & Dataset preparation and baseline retrieval pipeline \\
3--4 & CNN re-ranker implementation and training \\
5    & Experiments, ablations, and evaluation \\
6    & Report writing and final presentation \\
\bottomrule
\end{tabular}

\end{document}
